\section{طراحی آزمایش و آماده‌سازی داده}
\label{sec:corrected-protocol}
مبنای گزارش متمرکز، کد نوت‌بوک \en{finding-the-best-model-fixed.ipynb} است \cite{notebook-models}. مجموعه‌دادۀ \en{Suicide Watch} شامل ۲۳۲٬۰۷۴ ردیف خام با دو کلاس \en{non-suicide} و \en{suicide} است. برچسب‌ها به‌ترتیب به صفر و یک نگاشت می‌شوند. شمارش ردیف‌ها از گزارش بارگذاری، حذف تکرار و جداسازی بخش‌ها به دست می‌آید.

متن‌های تهی و برچسب‌های نامعتبر بررسی می‌شوند. برای تشخیص تکرار خام، نرمال‌سازی سازگاری یونیکد، یکسان‌سازی حالت حروف و فاصله‌ها انجام می‌شود. متن‌های تکراری با برچسب متعارض کنار گذاشته می‌شوند و از هر متن تکراریِ هم‌برچسب فقط یک ردیف باقی می‌ماند. ۹۳ ردیف تکراری خام حذف می‌شود و ۲۳۱٬۹۸۱ ردیف برای تقسیم اولیه باقی می‌ماند.

تقسیم نخست، ۲۰ درصد را برای آزمون نگه می‌دارد. از ۸۰ درصد توسعه، ۲۰ درصد برای اعتبارسنجی کنار گذاشته می‌شود؛ بنابراین سهم‌های اسمی آموزش، اعتبارسنجی و آزمون ۶۴، ۱۶ و ۲۰ درصد است. بذر تقسیم ۴۲ و روش تقسیم طبقه‌بندی‌شده است. پیش از پاکسازی، تعداد ردیف‌های این سه بخش به‌ترتیب ۱۴۸٬۴۶۷، ۳۷٬۱۱۷ و ۴۶٬۳۹۷ است. پس از پاکسازی و حذف هم‌پوشانی، تعدادهای نهایی به‌ترتیب ۱۴۸٬۳۸۰، ۳۷٬۰۸۶ و ۴۶٬۳۴۲ هستند. جدول~\ref{tab:fl-split} توزیع کلاس‌ها را نشان می‌دهد.

\subsection{پاکسازی ثابت و جلوگیری از نشت داده}
پاکسازی هر متن مستقل از برچسب و دیگر متن‌ها انجام می‌شود. ابتدا کدهای \en{HTML} بازگشایی و یونیکد با \en{NFKC} نرمال می‌شود. حروف به حالت کوچک تبدیل می‌شوند. نشانه‌های نقل‌قول یکسان‌سازی و برچسب‌های \en{HTML}، نشانی وب، نشانی رایانامه و شناسه‌های کاربری با فاصله جایگزین می‌شوند. صورت‌های متداول نفی در انگلیسی پیش از حذف نشانه‌گذاری باز می‌شوند؛ برای نمونه، \en{can't} به \en{can not} و \en{won't} به \en{will not} تبدیل می‌شود. سپس نویسه‌های غیرواژگانی حذف، زیرخط به فاصله تبدیل و فاصله‌های پیاپی یکسان می‌شوند.

اعداد و حروف یونیکد در پاکسازی حفظ می‌شوند و فهرست آموخته‌شده‌ای از کلمات توقف وجود ندارد. حذف بی‌قید واژۀ نفی می‌تواند معنای متن را وارونه کند؛ بنابراین فهرست توقف «دنا» در مسیر اصلی این آزمایش فراخوانی نمی‌شود. «دنا» یک ابزار جانبی فارسی است که در پیوست معرفی می‌شود و نتیجۀ آن نباید با نتیجۀ آزمایش انگلیسی حاضر یکی فرض شود.

پس از پاکسازی، ترتیب رسیدگی ثابتِ آموزش، اعتبارسنجی و آزمون اعمال می‌شود. متن‌های خالی و تکرارهای محلی حذف می‌شوند؛ اگر پاکسازی باعث یکسان شدن متنی با بخش قبلی شود، نمونه از بخش بعدی حذف می‌شود. متن‌های دارای برچسب متعارض تنها در همان بخش بررسی می‌شوند. بدین ترتیب، متن نگه‌داشته‌شدۀ آزمون سبب حذف یا تغییر نمونۀ آموزشی نمی‌شود. نبود هم‌پوشانی هم بر اساس شناسۀ ردیف و هم بر اساس متن پاکسازی‌شده کنترل می‌شود.

شکل~\ref{fig:cleaning} جریان تابع \en{clean\_text} و ممیزی بخش‌ها را نشان می‌دهد. اعداد و واژه‌های نفی در این مسیر حفظ می‌شوند.
\begin{figure}[htbp]
\centering\begin{tikzpicture}[
  node distance=5mm,
  stage/.style={draw=teal!60!black,fill=teal!3,rounded corners=3pt,
    inner sep=6pt,font=\small},
  flow/.style={-{Latex[length=2.2mm]},thick,draw=teal!60!black}]
\node[stage] (unicode) {\diagramtext{بازگشایی \en{HTML}، نرمال‌سازی \en{NFKC} و کوچک‌سازی حروف}};
\node[stage,below=of unicode] (links) {\diagramtext{حذف نشانی وب، رایانامه، برچسب \en{HTML} و شناسۀ کاربری}};
\node[stage,below=of links] (negation) {\diagramtext{باز کردن نفی انگلیسی پیش از حذف نشانه‌گذاری\\برای نمونه: \en{can't} به \en{can not}}};
\node[stage,below=of negation] (tokens) {\diagramtext{حفظ حروف یونیکد، اعداد و واژۀ هشتگ\\تبدیل نشانه‌گذاری و زیرخط به فاصله؛ یکسان‌سازی فاصله‌ها}};
\node[stage,below=of tokens] (audit) {\diagramtext{در هر بخش: حذف متن خالی، تعارض و تکرار محلی\\حذف هم‌پوشانی از بخش بعدی با اولویت آموزش، اعتبارسنجی، آزمون}};
\node[stage,below=of audit] (out) {\diagramtext{ممیزی جدایی شناسه و متن پاکسازی‌شده\\بدون فهرست توقف آموخته‌شده یا یادگیری واژگان از ارزیابی}};
\foreach \a/\b in {unicode/links,links/negation,negation/tokens,tokens/audit,audit/out}
  \draw[flow] (\a) -- (\b);
\end{tikzpicture}

\caption{جریان پاکسازی و ممیزی نسخۀ اصلاح‌شده؛ قواعد به برچسب ارزیابی وابسته نیستند}
\label{fig:cleaning}
\end{figure}

\subsection{واژه‌سازی و دامنۀ زبانی}
پس از پاکسازی، رشته بر اساس فاصله به توکن‌ها تقسیم می‌شود. این واژه‌ساز کم‌هزینه و قابل‌بازتولید است، اما تجزیۀ صرفی، رفع ابهام، تشخیص موجودیت یا مدل‌سازی زیرواژه‌ای انجام نمی‌دهد. خطاهای املایی و تفاوت صورت‌های واژگانی می‌توانند اندازۀ واژگان را افزایش دهند و واژه‌های کم‌فراوان را از پوشش مدل خارج کنند. یکسان‌سازی فاصله‌ها و نشانه‌ها این مسئله را کاهش می‌دهد، اما جایگزین تحلیل زبانی کامل نیست.

حفظ حروف یونیکد در کد به معنای ارزیابی موفق همۀ زبان‌ها نیست. دادۀ این آزمایش انگلیسی است و معیارهای آن نباید به فارسی، عبری یا عربی تعمیم داده شوند. برای چنین ادعایی به پیکرۀ مستقل، برچسب‌های معتبر و ارزیابی اختصاصی هر زبان نیاز است. این تمایز، مرز میان قابلیت فنی پردازش رشتۀ یونیکد و شاهد تجربی چندزبانه بودن را روشن می‌کند.

\section{بردارسازی مبتنی بر \en{Word2Vec}}
\label{sec:word2vec}
\en{Word2Vec} یک روش بازنمایی توزیعی واژه‌هاست که روابط هم‌وقوعی واژگان را از پیکره می‌آموزد \cite{ref78}. به هر واژۀ $w$ برداری $v_w\in\R^{d}$ اختصاص می‌یابد. در این پژوهش $d=100$ است. تمام طبقه‌بندهای این آزمایش ورودی صدبعدی دارند.

در معماری \en{CBOW}، واژۀ هدف از بافت پیش‌بینی می‌شود و در معماری \en{Skip-Gram}، واژه‌های بافت از واژۀ مرکزی پیش‌بینی می‌شوند. پیاده‌سازی حاضر از \en{Skip-Gram} با \en{sg=1} استفاده می‌کند. صورت مفهومی تابع هدف آن را می‌توان چنین نوشت:
\begin{equation}
\max\sum_{t=1}^{T}\sum_{\substack{-c\le j\le c\\j\ne0}}
\log P(w_{t+j}\mid w_t),\qquad
P(w_o\mid w_i)=\frac{\exp(u_{w_o}^{\mathsf T}v_{w_i})}
{\sum_{w\in V}\exp(u_w^{\mathsf T}v_{w_i})}.
\label{eq:skipgram}
\end{equation}
این رابطه توضیح مفهومی هدف توزیعی است؛ کتابخانه می‌تواند برای کاهش هزینه از تقریب‌هایی مانند نمونه‌گیری منفی استفاده کند. مدل حاضر شبکۀ کم‌عمقِ واژه‌محور است و همان مدل زمینه‌مند چندلایۀ خانوادۀ \en{Transformer} نیست.

پارامترهای بازنمایی عبارت‌اند از پنجرۀ ۵، حداقل فراوانی ۲، پنج دورۀ آموزش، یک کارگر و بذر ۴۲ برای مدل متمرکز. پنجره، محدودۀ بافتی را مشخص می‌کند؛ حداقل فراوانی، واژه‌های بسیار کم‌تکرار را از واژگان کنار می‌گذارد. افزایش ابعاد یا تعداد دوره‌ها می‌تواند هزینۀ محاسبات و حافظه را افزایش دهد و به‌تنهایی تضمین‌کنندۀ بهبود طبقه‌بندی نیست.

\subsection{تبدیل بردار واژه به بردار سند}
اگر سند $x$ شامل توکن‌های $w_1,\ldots,w_n$ باشد و $I(x)$ مجموعۀ جایگاه‌های دارای واژۀ شناخته‌شده در واژگان آموزشی $V$ باشد، نمایش سند برابر است با
\begin{equation}
\phi(x)=
\begin{cases}
\displaystyle\frac{1}{|I(x)|}\sum_{j\in I(x)}v_{w_j},&|I(x)|>0,\\[5pt]
\mathbf0\in\R^{100},&|I(x)|=0.
\end{cases}
\label{eq:document-embedding}
\end{equation}
مخرج، تعداد توکن‌های شناخته‌شده است؛ تقسیم بر تعداد تمام توکن‌ها با وجود واژه‌های ناشناخته، همان الگوریتم اجراشده نیست. توکن تکرارشده به تعداد وقوع خود در میانگین نقش دارد. طول سند تغییر می‌کند، اما طول بردار نهایی ثابت و برابر ۱۰۰ است.

میانگین‌گیری ترتیب واژه‌ها را در نمایش نهایی حفظ نمی‌کند. دو متن با واژگان مشابه و ترتیب متفاوت ممکن است بردارهای نزدیکی داشته باشند. چندمعنایی واژه‌ها نیز در بازنمایی ایستای واژه به‌صورت کامل حل نمی‌شود. در مقابل، هزینۀ محاسباتی و حافظه نسبت به بسیاری از مدل‌های جمله‌محور کمتر است و نمایش فشرده برای طبقه‌بندهای درختی و بیزی فراهم می‌شود. انتخاب این روش، یک مبادلۀ روش‌شناختی است و نباید برتری مطلق آن بر \en{TF-IDF} یا \en{Sentence-BERT} نتیجه گرفته شود.

\subsection{محل یادگیری بازنمایی و ممیزی واژگان}
توکن‌سازی، \en{MeanWord2Vec} و طبقه‌بند در یک شیء \en{Pipeline} قرار می‌گیرند. در هر تقسیم اعتبارسنجی متقابل، \en{fit} بازنمایی فقط بر متن‌های بخش آموزشی همان تقسیم اجرا می‌شود. متن‌های بخش اعتبارسنجی آن تقسیم صرفاً \en{transform} می‌شوند. پس از انتخاب فراپارامترها، مدل نهایی روی آموزش کامل بازبرازش می‌شود. حافظۀ نهان فقط برای بازاستفاده از تبدیل‌های یکسانِ همان تقسیم است و اجازۀ استفاده از واژگان بخش نگه‌داشته‌شده را نمی‌دهد.

در پایان، واژگان مدل منتخب با مجموعۀ توکن‌های آموزش مقایسه می‌شوند. باید $V\subseteq V_{\mathrm{train}}$ و $V\cap(V_{\mathrm{heldout}}\setminus V_{\mathrm{train}})=\varnothing$ برقرار باشد. پوشش توکن و سهم سندهای کاملاً ناشناخته جداگانه گزارش می‌شوند. این کنترل، یک آزمون مشخص برای نشت واژگان است؛ جایگزین تمام بررسی‌های ممکن نشت داده یا تعمیم‌پذیری نیست.

\section{معیارهای ارزیابی و انتخاب مدل}
\label{sec:metrics}
کلاس مثبت، برچسب \en{suicide} با مقدار یک است. اگر $TP$ و $TN$ تعداد پیش‌بینی‌های صحیح مثبت و منفی و $FP$ و $FN$ تعداد مثبت و منفی کاذب باشند، معیارهای اصلی از روابط زیر محاسبه می‌شوند:
\begin{align}
\mathrm{Accuracy}&=\frac{TP+TN}{TP+TN+FP+FN},\\
\mathrm{Precision}_{+}&=\frac{TP}{TP+FP},&
\mathrm{Recall}_{+}&=\frac{TP}{TP+FN},\\
F_{1,+}&=\frac{2TP}{2TP+FP+FN},&
\mathrm{Specificity}&=\frac{TN}{TN+FP}.
\label{eq:classification-metrics}
\end{align}
در حالت مخرج صفر، مطابق پیاده‌سازی، مقدار صفر گزارش می‌شود. میانگین کلان، میانگین سادۀ معیار دو کلاس است؛ میانگین وزن‌دار از تعداد نمونۀ واقعی هر کلاس وزن می‌گیرد. دقت متوازن نیز میانگین بازیابی دو کلاس است. در دادۀ نامتوازن، دقت کلی به‌تنهایی می‌تواند گمراه‌کننده باشد؛ بنابراین بازیابی و $F_1$ کلاس مثبت، تعداد خطاها و معیارهای رتبه‌بندی نیز بررسی می‌شوند \cite{ref111,ref117,ref118,ref119}.

\en{ROC AUC} توان رتبه‌بندی نمونۀ مثبت بالاتر از نمونۀ منفی را در دامنه‌ای از آستانه‌ها خلاصه می‌کند. \en{Average Precision} یا \en{AP} خلاصه‌ای از منحنی دقت مثبت\lr{--}بازیابی است؛ با مساحت ذوزنقه‌ای هر منحنی دلخواه یکسان فرض نمی‌شود. این معیارها از احتمال‌ها محاسبه می‌شوند، در حالی که دقت، بازیابی، $F_1$ و ماتریس سردرگمی به آستانه وابسته‌اند.

\subsection{سه سطح جداسازی ارزیابی}
فراپارامترهای درخت تصمیم و \en{XGBoost} با جست‌وجوی تصادفی ۱۰ ترکیب و اعتبارسنجی طبقه‌بندی‌شدۀ سه‌بخشی فقط روی آموزش تنظیم می‌شوند. معیار این جست‌وجو، $F_1$ کلاس مثبت در آستانۀ پیش‌فرض طبقه‌بند است. سپس برای هر مدلِ بازبرازش‌شده، آستانه از ۱۰۱ مقدار صفر تا یک روی اعتبارسنجی مستقل انتخاب می‌شود. در آخر، مدل‌ها با ترتیب $F_{1,+}$، $F_{1,\mathrm{macro}}$، \en{ROC AUC} و نام مدل رتبه‌بندی می‌شوند. انتخاب نهایی پیش از مشاهدۀ معیار آزمون تثبیت می‌شود \cite{ref122}.

دادۀ آزمون برای انتخاب بهترین مدل از میان چهار مدل استفاده نمی‌شود. به همین دلیل، جدول مقایسۀ هر چهار مدل در این نسخه مربوط به اعتبارسنجی است و ارزیابی نهایی آزمون برای مدل منتخب ارائه می‌شود. مقایسۀ مدل‌ها باید بر مجموعۀ مشترک و با نقش ارزیابی یکسان انجام شود.
